% \appendix
% \section{appendix}
\section{Model Details}


Acoustic Encoder: Following the design principles in \cite{kumar2024high}, our encoder comprises four convolutional encoder blocks, each tailored to progressively downsample the input audio waveform at rates [2, 4, 5, 8]. This structured reduction ensures efficient encoding while preserving essential audio characteristics. The final output from the encoder has a hidden size of 256, ensuring detailed feature representation. The total model size for the encoder stands at 12.18MB.

Residual Vector Quantizer (RVQ): A key component in our X-codec, the RVQ utilizes the techniques established by \cite{zeghidour2021soundstream}. We update the codebook entries using an exponential moving average and apply a straight-through estimator to facilitate the gradient computation during backpropagation. To bolster training effectiveness and adaptability, commitment loss is incorporated, and RVQ layers are randomly selected from options [1, 2, 3, 4, 8] during training. This variability allows the model to adapt more dynamically to different audio characteristics.

Acoustic Decoder: The decoder is designed to mirror the encoder’s architecture, featuring four layers that upsample the audio data at inverse rates [8, 5, 4, 2]. This symmetry between the encoder and decoder helps in effectively reconstructing the audio signal from its encoded state. The decoder's model size is approximately 19.27 MB.

Semantic Encoder and Decoder: To further refine the semantic aspects of the audio signals, we incorporate two additional convolutional blocks within both the semantic encoder and decoder, each with a hidden size of 768. This setup enhances the model’s ability to process and integrate semantic information effectively.
\section{Music Continue with 4 VQ Flattern}
% xcodec 
% frechet_audio_distance: 0.2029561
% kullback_leibler_divergence_sigmoid: 0.7827974
% kullback_leibler_divergence_softmax: 0.1475402
% inception_score_mean: 1.5393551
% inception_score_std: 0.0818533
% frechet_distance: 3.4788028

% 0.7972045598772866 

% frechet_audio_distance: 7.0863215
% kullback_leibler_divergence_sigmoid: 0.8706447
% kullback_leibler_divergence_softmax: 0.1659164
% inception_score_mean: 1.6810636
% inception_score_std: 0.1291735
% frechet_distance: 3.8605033

% 1.1233672557857517   
In this section, we expand our music continuation experiments to include conditions with four flattened vector quantizations (VQs) to further validate the effectiveness of our approach. While the experimental details remain consistent with previous setups, the use of four VQs necessitates a shorter segment length of approximately 20 seconds due to increased data density. We prompt the model with 5 seconds of audio and generate  5 seconds, aiming to assess the performance enhancement under these conditions.

\begin{table}[h]
\centering
\caption{Comparison between baseline acoustic codec and our X-codec on music continuation.}
\resizebox{0.4\textwidth}{!}{%
\begin{tabular}{lccc}
\hline
\textbf{Model} & \textbf{FD} $\downarrow$ & \textbf{FAD} $\downarrow$ & \textbf{FD-MERT-layer-9}  $\downarrow$  \\
\hline
Acoustic codec   &3.86  & 7.08 & 1.12  \\
X-codec & 3.47 & 0.20  &0.79   \\
\hline
\end{tabular}
}
\label{music}
\end{table}
The results underscore a noticeable improvement in performance with the X-codec, particularly highlighted by significant enhancements in Frechet Audio Distance (FAD) and perceptual quality metrics, suggesting that our X-codec not only maintains but also amplifies its efficacy in generating musically coherent and contextually rich outputs.
